\chapter{Weather and Climate}
\section{Get Rolling}
\textbf{Discuss the following questions in the classroom.}
\begin{enumerate}
    \item Reema is feeling cold and wearing warm clothes. What season is it?\\
    Answer: Winter
    \item The sun is shining bright in the sky. What weather is it?\\
    Answer: Sunny
    \item There are dark clouds in the sky and the Sun is not visible. What weather is it?\\
    Answer: Cloudy
\end{enumerate}
\section{What to Know}
Earth's climate is influenced by different factors. Earth can be divided into various climatic zones based on how much heat each region recives from the Sun.
\section{\textbf{Learning Objectives}}
\begin{itemize}
    \item to differentiate between climate and weather
    \item to explain the main factors that affect climate
    \item to discuss the three climatic zones on Earth
    \item to find latitudes of any given place on a globe
    \item to analyse the impact of climate on people
\end{itemize}
\section{I Explore}
\section{Weather}
\textbf{Weather}
is defined as the condition of the atmosphere at a particular place at a particular time.
\vspace{0.1in}
\par
There are days when it is hot and sunny, and there are days when it suddenly starts raining and becomes cool. On some days, a dust storm in the morning may be followed by clear skies in the evening. This is because the weather of a place keeps changing every day. Sometimes, it even changes several times in the same day. Weather may change from place to place, time to time, day to day or season to season.
\vspace{0.1in}
\par
\textbf{Meteorologists} keep daily records of temperatures, precipitation and the wind. From this data, they get an idea of the weather.
\section{Climate}
\textbf{Climate} refers to the average weather conditions of a place that continue for a long time, usually 30 years. Different parts of the world have different types of climates. Temperature, humidity, precipitation, wind speed and direction are some of the elements of weather and climate.
\subsection{Factors Influencing Climate}
The climate differs from place to place. This is because many factors influence the distribution of heat on Earth.
\subsection{\textbf{Distance from the Equator or latitude}}
The rays of the Sun fall directly on the places situated near the Equator. Since these rays are focused over a small area, it gets most of the Sun's heat and light. On the other hand, when sunrays strike areas further away from the Equator, they fall slanting and spread out over a larger area. Hence, their heating effect is less. Therefore, places that are far away from the Equator are cooler than those near the Equator.
\subsection{\textbf{Distance from the sea}}
Seas and Oceans affect the climate of places close them. Sea water heats up and cools down slower than land. During the night, the sea is warmer than the land. The air above the sea gets warm, becomes lighter and rises. Cool air from the land flows to fill the empty space. As the air moves from the land to the sea during night, it is called \textbf{land breeze.} During the day, the land is warmer than the sea. The air above the land gets warm, becomes lighter and rises. Cool air from the sea is heavier because it has moisture. It blows from the sea to the land and fills the empty space. As the air moves from the sea to the land during the day, it is called \textbf{sea breeze.} Therefore, costal regions or places close to the sea have a moderate climate. They are neither too hot in summers nor too cold in winters. However, places far away from the seas and oceans are landlocked areas, and thus they do not experience the moderating influence of seas and oceans. They experience an \textbf{extreme climate,} with very hot summers and extremely cold winters.
\vspace{0.1in}
\par
I wonder what clothes we should carry while travelling to Mumbai and Chandigarh during winters!
\subsection{\textbf{Height above the sea level}}
Temperature decreases with increasing altitude or the height above sea level. It drops by 1 degree Celsius for every 100 meters of altitude. Higher places, or places at higher altitudes, are usually cooler. Therefore, mountains are colder than plains. The climate of a place is affected not only by latitudes but also by altitude. For example, though Mt Kilimanjaro in Africa is located near the Equator, its peak is covered in snow. Hill stations like Naintal, Darjeeling, Ooty and Shimla are cool even in summers.
\subsection{\textbf{Direction of winds}}
The winds blowing from cold regions bring cool air. The winds coming from hot region bring hot air. For example, cold winds from Himalayas make northern India very cold. Hot winds from Rajasthan make Delhi hot in summers. The winds blowing from the sea or ocean make the places near it cool and humid. Sometimes the winds bring clouds with them and cause rainfall.
\subsection{\textbf{Humidity and rainfall}}
\textbf{Humidity} refers to the presence of water vapour in the air. Humidity in the air depends on the location of a place.
\vspace{0.1in}
\par
Places near the Equator are hot and humid and receive a lot of rainfall throughout the year. Air over the deserts do not have moisture, so these places receive very little or no rainfall. A place located near large water bodies, like oceans, seas and rivers, would be more humid than places situated far away from them.
\vspace{0.4in}
\par
Humidity also depends on the direction of wind. The winds that blow from the sea to land usually bring a lot of moisture with them and cause rainfall, whereas the winds that blow from land to the sea are usually dry and do not cause any rain.
\section{Heat Zones of the World}
Earth can be divided into three zones based on the amount of heat a region gets from the Sun. These zones are called \textbf{heat zones} or \textbf{climatic zones.}
\subsection{Torrid Zone}
The torrid zone or the tropical zone is located between the Tropic of Cancer and the Tropic ofcCapricorn. It lies on both sides of the Equator. The climate of this region is the hottest because it receives direct and vertical rays of the Sun.
\subsection{Temperate Zone}
The temperate zone is the area located between the Tropic of Cancer and the Arctic Circle in the Northern Hemisphere and between the Tropic of Capricorn and the Antarctic Circle in the Southern Hemisphere. This zone receives slanting rays of the Sun. Hence, this region experiences a moderate climate-neither too hot nor too cold.
\subsection{Frigid Zone}
The frigid zone or the polar zone is the region that is located between the Arctic Circle and the North Pole in the Northern Hemisphere and between the Antarctic Circle and the South Pole in the Southern Hemisphere. Thus, there is a north frigid zone and a south frigid zone.
\vspace{0.4in}
\par
Being the farthest from the Equator, these areas recive slanting rays of the Sun. That is why the frigid zones are the coldest parts of the Earth and remain covered in snow and ice.
\section{Impact of Climate and Season}
Changes in the climate and season influence the lifestyle of people. The food habits, clothing styles, houses and activities of people are affected by the climate of the place they live in.
\vspace{0.1in}
\par
Winters are the time to enjoy hot foods with family and friends. Summer hat sweetens mangoes and watermelons. Farmers grow crops according to the seasons.
\vspace{5.6in}
\section{Now I Know}
\begin{itemize}
    \item Weather is defined as the condition of the atmosphere at a particular place at a particular time.
    \item Climate is the average or general weather condition of a place that continues for a long time, usually 30 years. Climate does not change as frequently as weather.
    \item The Climate of a place depends upon many factors-its distance from the Equator, its distance from the sea, its height above sea level, direction of winds, and humidity and rainfall.
    \item Earth can be divided into three zones based on the amount of heat a region gets from the Sun. These zones are called the heat zones or the climatic zones. The three heat zones are the torrid zone, the temperate zone and the frigid zone.
    \item Changes in weather and seasons have an impact on the habits and living conditions of people.
\end{itemize}